\paragraph{分歧点的有理降阶、斜率平方不变量与 Lee--Yang 边界的普适定标核}
本段以一参数变形谱四次
\[
F_t(\mu,u):=\mu^4-\mu^3-(tu+1)\mu^2+\mu(tu-u^3)+tu
\]
及其分歧系统
\(
F_t(\mu,u)=\partial_\mu F_t(\mu,u)=0
\)
为背景记号，抽取两类内禀对象：由双根约束诱导的有理消元律与平面代数曲线化表述，以及由 Puiseux 分歧诱导的 Lee--Yang 斜率平方不变量 \(b_\ast\)。
在主导并合分支点处，它们进一步给出 \(\widetilde M_m(u)\) 的两类普适定标核与二次晶格型零点律。
\par
在折叠语义的 Bernoulli 半 Perron 谱模型中，沿用 \eqref{eq:fold-bernoulli-half-perron-quartic} 的记号并取特化 \(t=2\)，记
\(
F(\mu,u):=F_2(\mu,u)
\)；
本段后续关于 \(Q_{10}\)、\(P_{10}\) 与 \(b_\ast\) 的对象均在该特化下讨论。

\begin{proposition}[双根约束的线性消元闭式与降阶 Tschirnhaus 同余]\label{prop:fold-gauge-anomaly-branchpoint-rational-reduction}
设 \(u\neq 0\)。若存在 \(\mu\in\CC\) 使得 \(\mu\) 为 \(F_t(\cdot,u)\) 的重根，即
\[
F_t(\mu,u)=0,\qquad \partial_\mu F_t(\mu,u)=0,
\]
则必有 \(\mu\neq \pm i\)。令
\[
X:=tu,\qquad Y:=u^3,
\]
则 \((X,Y)\) 由 \(\mu\) 唯一确定，并满足显式有理公式
\[
\boxed{\;
X=\frac{\mu^2(3\mu^2-2\mu-1)}{\mu^2+1},\qquad
Y=-\frac{2\mu(\mu^2-\mu-1)^2}{\mu^2+1}\; }.
\]
反之，给定任意 \(\mu\neq \pm i\)，若存在 \((t,u)\) 使得 \(tu=X\) 且 \(u^3=Y\)，则 \(\mu\) 必为 \(F_t(\cdot,u)\) 的重根。
\par
在特化 \(t=2\) 且 \(\mu\neq \pm i\) 时，由 \(u=X/t\) 得到降阶有理消元律
\begin{equation}\label{eq:fold-gauge-anomaly-u-rational-reduction}
\boxed{\;
u=\mathcal R(\mu):=\frac{\mu^2(3\mu^2-2\mu-1)}{2(\mu^2+1)}\; }.
\end{equation}
进一步，令 \(Q_{10}(\mu)\) 与 \(P_{10}(u)\) 如命题 \ref{prop:fold-gauge-anomaly-Q10-tschirnhaus}。
则对任意根 \(\mu_\ast\) 满足 \(Q_{10}(\mu_\ast)=0\)，都有 \(u_\ast=\mathcal R(\mu_\ast)\) 满足 \(P_{10}(u_\ast)=0\)；
并且 \(\mathcal R\) 与命题 \ref{prop:fold-gauge-anomaly-Q10-tschirnhaus} 中的 Tschirnhaus 变换 \(T(\mu)\) 在 \(Q_{10}=0\) 的根集上严格一致：存在整系数恒等式
\begin{equation}\label{eq:fold-gauge-anomaly-R-minus-T-congruence}
\boxed{\;
\mathcal R(\mu)-T(\mu)=\frac{(7\mu+1)\,Q_{10}(\mu)}{200(\mu^2+1)}\; }.
\end{equation}
\end{proposition}
\leanverified{paper\_fold\_gauge\_anomaly\_branchpoint\_rational\_reduction}
\begin{proof}
将 \(F_t(\mu,u)=0\) 与 \(\partial_\mu F_t(\mu,u)=0\) 中出现的 \(tu\) 与 \(u^3\) 记为未知量 \(X,Y\)。
两式等价改写为关于 \((X,Y)\) 的二元一次方程组
\[
\mu^4-\mu^3-\mu^2-\mu Y+(-\mu^2+\mu+1)X=0,
\qquad
4\mu^3-3\mu^2-2\mu-Y+(1-2\mu)X=0.
\]
其系数行列式为 \(\mu^2+1\)，故当且仅当 \(\mu\neq \pm i\) 时该方程组有唯一解。
直接消元得到所示 \(X(\mu),Y(\mu)\) 的闭式表达；再将解代回方程组即得“反向”结论。
\par
对 \eqref{eq:fold-gauge-anomaly-u-rational-reduction}，取特化 \(t=2\) 并由 \(u=X/t\) 立即得到。
\par
对 \eqref{eq:fold-gauge-anomaly-R-minus-T-congruence}，将 \(T(\mu)\) 取为命题 \ref{prop:fold-gauge-anomaly-Q10-tschirnhaus} 的显式九次多项式除以 \(200\)，通分并化简即可得到所示整系数恒等式；
右端显式含因子 \(Q_{10}(\mu)\)，故在 \(Q_{10}(\mu)=0\) 的根集上两者一致。
从而对任意 \(Q_{10}\)-根 \(\mu_\ast\)，有 \(\mathcal R(\mu_\ast)=T(\mu_\ast)\)，并由命题 \ref{prop:fold-gauge-anomaly-Q10-tschirnhaus} 得 \(P_{10}(\mathcal R(\mu_\ast))=0\)。证毕。
\end{proof}

\begin{corollary}[\texorpdfstring{$\mu=-1$}{mu=-1} 分歧共振的全复域分类与实刚性]\label{cor:fold-gauge-anomaly-mu-minus1-branch-classification}
在命题 \ref{prop:fold-gauge-anomaly-branchpoint-rational-reduction} 的设定下，若重根取 \(\mu=-1\)，则双根条件等价于
\[
tu=2,\qquad u^3=1.
\]
因此 \(\mu=-1\) 的重根事件在复参数域中恰好发生于三点族
\[
\boxed{\ (t,u)=\bigl(2/\zeta,\ \zeta\bigr),\qquad \zeta^3=1\ },
\]
并且在实参数截面上仅有 \((t,u)=(2,1)\)。
\end{corollary}
\begin{proof}
由命题 \ref{prop:fold-gauge-anomaly-branchpoint-rational-reduction} 的显式公式直接代入 \(\mu=-1\) 得 \(X=2\)、\(Y=1\)，等价于 \(tu=2\)、\(u^3=1\)。
反向由定义代回可知此时 \(\mu=-1\) 同时使 \(F_t\) 与 \(\partial_\mu F_t\) 为零。
实参数情形 \(u^3=1\Rightarrow u=1\) 从而 \(t=2\)。证毕。
\end{proof}
\leanverified{paper\_fold\_gauge\_anomaly\_mu\_minus1\_branch\_classification}

\begin{proposition}[\texorpdfstring{$u$}{u}-判别式的 \texorpdfstring{$\mu$}{mu}-降维与有理回收]\label{prop:fold-gauge-anomaly-branchpoint-mu-reduction-Et}
\leanverified{paper\_fold\_gauge\_anomaly\_branchpoint\_mu\_reduction\_et}
设 \(t\neq 0\)。则 \(u\) 为多项式 \(F_t(\cdot,u)\) 在 \(\mu\)-方向出现重根的分歧点当且仅当存在 \(\mu\neq \pm i\) 满足单一方程
\[
\boxed{\;
E_t(\mu):=
2t^3(\mu^2+1)^2(\mu^2-\mu-1)^2+\mu^5(3\mu^2-2\mu-1)^3=0\; }.
\]
在该情形下分歧点可由 \(\mu\) 有理回收为
\[
\boxed{\;
u=\frac{\mu^2(3\mu^2-2\mu-1)}{t(\mu^2+1)}\; }.
\]
特别地，当 \(t=2\) 时，
\[
E_2(\mu)=(\mu+1)\,Q_{10}(\mu),
\]
其中 \(Q_{10}\) 与命题 \ref{prop:fold-gauge-anomaly-Q10-tschirnhaus} 的定义一致，并且 \(\mu=-1\) 对应分歧点 \(u=1\)。
\end{proposition}
\begin{proof}
由命题 \ref{prop:fold-gauge-anomaly-branchpoint-rational-reduction}，在 \(\mu\neq \pm i\) 下双根条件等价于 \(X=tu\) 与 \(Y=u^3\) 取为所给有理表达。
由 \(t\neq 0\) 得
\(
u=X/t=\mu^2(3\mu^2-2\mu-1)/(t(\mu^2+1))
\)。
将其立方并与 \(Y\) 相等，清分母即得 \(E_t(\mu)=0\)。
反向：若 \(E_t(\mu)=0\)，以上定义的 \(u\) 满足 \(u^3=Y\)，从而 \((X,Y)\) 同时满足二元一次方程组，等价于 \(F_t(\mu,u)=\partial_\mu F_t(\mu,u)=0\)。
\par
当 \(t=2\) 时，对 \(E_2(\mu)\) 作整式因式分解得到 \((\mu+1)Q_{10}(\mu)\)；并由 \(\mu=-1\) 时的有理回收立得 \(u=1\)。证毕。
\end{proof}

\begin{proposition}[分歧曲线在单项式坐标 \texorpdfstring{$(X,Y)=(tu,u^3)$}{(X,Y)=(tu,u^3)} 下的显式代数方程与有理性]\label{prop:fold-gauge-anomaly-branch-curve-XY-explicit}
\leanverified{paper\_fold\_gauge\_anomaly\_branch\_curve\_xy\_explicit}
令 \(X=tu\)、\(Y=u^3\)。则 \(F_t(\cdot,u)\) 在 \(\mu\)-方向出现重根的充要条件等价于平面代数曲线
\[
R(X,Y)=0,
\]
其中 \(R\in\ZZ[X,Y]\) 可取为如下显式不可约多项式：
\[
\begin{aligned}
R(X,Y)=\;&20X^5-8X^4Y-120X^4+4X^3Y^2+236X^3Y+220X^3\\
&-239X^2Y^2-164X^2Y-120X^2+90XY^3-70XY^2-108XY+20X\\
&-27Y^4-22Y^3+5Y^2.
\end{aligned}
\]
并且该曲线为有理曲线：由 \(\mu\neq \pm i\) 给出有理参数化
\[
X(\mu)=\frac{\mu^2(3\mu^2-2\mu-1)}{\mu^2+1},\qquad
Y(\mu)=-\frac{2\mu(\mu^2-\mu-1)^2}{\mu^2+1}.
\]
\end{proposition}
\begin{proof}
命题 \ref{prop:fold-gauge-anomaly-branchpoint-rational-reduction} 给出 \((X,Y)\) 关于 \(\mu\) 的有理表达；消去参数 \(\mu\)（例如取两式的结果式）得到某个非零整系数消元多项式 \(R\) 使 \(R(X(\mu),Y(\mu))\equiv 0\)。
对所给显式 \(R\) 直接检验该恒等式成立。
由于给出了有理参数化，曲线函数域同构于 \(\CC(\mu)\)，从而曲线为有理曲线；不可约性可由参数化的函数域同构立即推出。证毕。
\end{proof}

\leanverified{paper\_fold\_gauge\_anomaly\_resonant\_branchpoint\_series}
\begin{proposition}[\texorpdfstring{$(t,u)=(2,1)$}{(t,u)=(2,1)} 共振点附近的解析分歧支与三阶展开]\label{prop:fold-gauge-anomaly-resonant-branchpoint-series}
考虑推论 \ref{cor:fold-gauge-anomaly-mu-minus1-branch-classification} 的实共振点 \((t,u)=(2,1)\)。
则存在唯一解析函数 \(u_b(t)\)（定义在 \(t=2\) 的充分小邻域内）满足 \(u_b(2)=1\) 且 \(F_t(\cdot,u_b(t))\) 在 \(\mu\)-方向出现重根；
相应的重根位置 \(\mu_b(t)\) 亦解析，且局部展开满足
\[
\boxed{\;
u_b(t)=1+(t-2)-\frac12(t-2)^2+\frac78(t-2)^3+O\!\bigl((t-2)^4\bigr)\;},
\]
\[
\boxed{\;
\mu_b(t)=-1-\frac12(t-2)+\frac{3}{16}(t-2)^2-\frac{23}{64}(t-2)^3+O\!\bigl((t-2)^4\bigr)\; }.
\]
从而 \(u_b'(2)=1\)。
\end{proposition}
\begin{proof}[证明要点]
由命题 \ref{prop:fold-gauge-anomaly-branchpoint-rational-reduction}，分歧曲线可在 \(\mu=-1\) 的邻域用参数 \(\mu\) 表示为 \(X(\mu)=tu\)、\(Y(\mu)=u^3\)。
在 \(Y(-1)=1\) 的邻域取满足 \(u(-1)=1\) 的解析分支 \(u(\mu)=Y(\mu)^{1/3}\)，并令 \(t(\mu)=X(\mu)/u(\mu)\)；
则 \((t(\mu),u(\mu))\) 参数化了穿过 \((2,1)\) 的分歧曲线局部支。
对 \(t(\mu)\)、\(u(\mu)\) 在 \(\mu=-1\) 处作泰勒展开并反演级数得到所示三阶展开；\(\mu_b(t)\) 由同一反演过程给出。证毕。
\end{proof}

\leanverified{paper\_fold\_gauge\_anomaly\_leyang\_slope\_square\_minpoly}
\begin{proposition}[Lee--Yang 斜率平方不变量与十次整系数极小多项式]\label{prop:fold-gauge-anomaly-leyang-slope-square-minpoly}
令 \((\mu_\ast,u_\ast)\) 为任一非平凡分歧点，满足
\[
u_\ast\notin\{0,1\},\qquad
F(\mu_\ast,u_\ast)=\partial_\mu F(\mu_\ast,u_\ast)=0,
\qquad
\partial_uF(\mu_\ast,u_\ast)\neq 0,\ \partial_{\mu\mu}^2F(\mu_\ast,u_\ast)\neq 0.
\]
定义 Puiseux 分歧诱导的斜率平方不变量
\begin{equation}\label{eq:fold-gauge-anomaly-bstar-definition}
\boxed{\;
b_\ast:=
\frac{-2\,\partial_u F(\mu_\ast,u_\ast)}{\partial_{\mu\mu}^2F(\mu_\ast,u_\ast)\,\mu_\ast^{2}}\; }.
\end{equation}
则 \(b_\ast\) 为下列十次整系数多项式的根：
\begin{equation}\label{eq:fold-gauge-anomaly-bstar-minpoly}
\boxed{\;
\begin{aligned}
0=\;&316283904\, b^{10}-52713984\, b^{9}+379196672\, b^{8}+406374016\, b^{7}\\
&+670701296\, b^{6}+1305085464\, b^{5}+761320448\, b^{4}+12917152\, b^{3}\\
&+53372880\, b^{2}+7820550\, b-1303425.
\end{aligned}}
\end{equation}
\end{proposition}
\begin{proof}
由命题 \ref{prop:fold-gauge-anomaly-branchpoint-rational-reduction}，任一非平凡分歧点满足 \(u_\ast=\mathcal R(\mu_\ast)\)，故 \eqref{eq:fold-gauge-anomaly-bstar-definition} 的 \(b_\ast\) 可化为 \(\mu_\ast\) 的有理函数值。
另一方面，\(\mu_\ast\) 必满足命题 \ref{prop:fold-gauge-anomaly-Q10-tschirnhaus} 的约束 \(Q_{10}(\mu_\ast)=0\)。
将方程 \(b-\mathcal B(\mu)=0\)（其中 \(\mathcal B\) 为上述有理表达）与 \(Q_{10}(\mu)=0\) 对 \(\mu\) 作结果式消元，得到仅含 \(b\) 的十次整系数多项式；化简后即为 \eqref{eq:fold-gauge-anomaly-bstar-minpoly}。证毕。
\end{proof}

\leanverified{paper\_fold\_gauge\_anomaly\_leyang\_universal\_kernels}
\begin{theorem}[Lee--Yang 边界的两类普适定标核：\texorpdfstring{$\sinh(w)/w$}{sinh(w)/w} 与 \texorpdfstring{$\cosh(w)$}{cosh(w)}]\label{thm:fold-gauge-anomaly-leyang-universal-kernels}
沿用命题 \ref{prop:fold-gauge-anomaly-mgf-order4-recurrence} 的记号，令 \(\widetilde M_m(u):=2^mM_m(u)\)。
取一非平凡分歧点 \((\mu_\ast,u_\ast)\) 位于主导谱支并合处，即在 \(u=u_\ast\) 时四根中存在两根作并合且其模长达到最大值 \(|\mu_\ast|\)（例如命题 \ref{prop:fold-gauge-anomaly-branch-point-classification} 的 Lee--Yang edge 共轭对）。
令 \(b_\ast\) 为 \eqref{eq:fold-gauge-anomaly-bstar-definition}。
则存在常数 \(C_1,C_0\in\CC\)（仅依赖于 \(\widetilde M_m\) 的初值与端点泛函，而与 \(m\) 无关），使得对任意紧集 \(K\subset\CC\)，当 \(m\to\infty\) 时在 \(w\in K\) 上一致成立以下二择一结论：
\begin{enumerate}
  \item 若 \(\displaystyle \lim_{m\to\infty}\frac{\widetilde M_m(u_\ast)}{m\,\mu_\ast^{m}}=C_1\neq 0\)，则
  \begin{equation}\label{eq:fold-gauge-anomaly-leyang-sinh-kernel}
  \boxed{\;
  \lim_{m\to\infty}\frac{\widetilde M_m\!\left(u_\ast+\frac{w^2}{b_\ast m^2}\right)}{m\,\mu_\ast^{m}}
  =2C_1\,\frac{\sinh(w)}{w}\; }.
  \end{equation}
  \item 若 \(\displaystyle \lim_{m\to\infty}\frac{\widetilde M_m(u_\ast)}{\mu_\ast^{m}}=C_0\neq 0\)，则
  \begin{equation}\label{eq:fold-gauge-anomaly-leyang-cosh-kernel}
  \boxed{\;
  \lim_{m\to\infty}\frac{\widetilde M_m\!\left(u_\ast+\frac{w^2}{b_\ast m^2}\right)}{\mu_\ast^{m}}
  =2C_0\,\cosh(w)\; }.
  \end{equation}
\end{enumerate}
\end{theorem}
\begin{proof}[证明要点]
在 \(u_\ast\) 的去心邻域内取局部参数 \(s=\sqrt{u-u_\ast}\)，由分歧非退化性得到两条主导谱支的 Puiseux 展开
\(
\mu_\pm(s)=\mu_\ast\exp(\pm \sqrt{b_\ast}\,s+O(s^2))
\)。
由四阶递推的谱分解，\(\widetilde M_m(u)\) 在该邻域可写为四个谱支贡献之和；其中两项为 \(\mu_\pm(s)^m\)，其余两项在 \(u_\ast\) 处模长严格小于 \(|\mu_\ast|\)，从而在 \(m\to\infty\) 的主导定标下指数可忽略。
由于 \(\widetilde M_m(u(s))\) 为 \(s\) 的偶函数，其在 \(\mu_\pm(s)^m\) 的系数允许出现至多一阶极点，故可整理为
\[
\widetilde M_m(u(s))
=a_{-1}(s^2)\frac{\mu_+(s)^m-\mu_-(s)^m}{s}
+a_0(s^2)\bigl(\mu_+(s)^m+\mu_-(s)^m\bigr)
+O(\rho^m),
\]
其中 \(|\rho|<|\mu_\ast|\) 且 \(a_{-1},a_0\) 在 \(s^2\) 上解析。
令 \(w=m\sqrt{b_\ast}\,s\)（等价于 \(u=u_\ast+w^2/(b_\ast m^2)\)），并用指数近似
\(
\mu_\pm(s)^m=\mu_\ast^m e^{\pm w}(1+o(1))
\)
即可得到 \eqref{eq:fold-gauge-anomaly-leyang-sinh-kernel} 与 \eqref{eq:fold-gauge-anomaly-leyang-cosh-kernel} 的两种主导情形。证毕。
\end{proof}

\begin{corollary}[Lee--Yang 零点的二次晶格律与方向不变量]\label{cor:fold-gauge-anomaly-leyang-quadratic-lattice-zeros}
在定理 \ref{thm:fold-gauge-anomaly-leyang-universal-kernels} 的设定下，存在 \(\delta_0>0\) 使得当 \(m\) 足够大时，\(\widetilde M_m(u)\) 在圆盘 \(D(u_\ast,\delta_0)\) 内的零点满足二次晶格定律：
\begin{enumerate}
  \item 在 \eqref{eq:fold-gauge-anomaly-leyang-sinh-kernel} 的情形下，对每个固定 \(k\in\NN\) 存在零点 \(u_{m,k}\in D(u_\ast,\delta_0)\) 使
  \[
  u_{m,k}=u_\ast-\frac{\pi^2 k^2}{b_\ast m^2}+O\!\left(\frac{k^3}{m^3}\right)\qquad(m\to\infty).
  \]
  \item 在 \eqref{eq:fold-gauge-anomaly-leyang-cosh-kernel} 的情形下，对每个固定 \(k\in\NN\) 存在零点 \(u_{m,k}\in D(u_\ast,\delta_0)\) 使
  \[
  u_{m,k}=u_\ast-\frac{\pi^2 (k+\tfrac12)^2}{b_\ast m^2}+O\!\left(\frac{k^3}{m^3}\right)\qquad(m\to\infty).
  \]
\end{enumerate}
因此零点以 \(m^{-2}\) 速度逼近 \(u_\ast\)，其逼近方向由复向量
\[
\boxed{\ \mathrm{dir}(u_\ast)=-\frac{1}{b_\ast}\ }
\]
完全决定。
\leanverified{paper\_fold\_gauge\_anomaly\_leyang\_quadratic\_lattice\_zeros}
\end{corollary}
\begin{proof}
定理 \ref{thm:fold-gauge-anomaly-leyang-universal-kernels} 给出经归一化后的解析函数在紧集上一致收敛到
\(
2C_1\sinh(w)/w
\)
或
\(
2C_0\cosh(w)
\)。
两者零点分别为 \(w=i\pi k\ (k\in\ZZ\setminus\{0\})\) 与 \(w=i\pi(k+\tfrac12)\ (k\in\ZZ)\)。
由一致收敛与 Rouch\'e 定理，极限零点在 \(m\) 足够大时稳定地产生对应的 \(w_{m,k}\)，再由解析反变换 \(u=u_\ast+w^2/(b_\ast m^2)\) 得到所示渐近式。证毕。
\end{proof}

\begin{corollary}[数域同构：\texorpdfstring{$u_\ast,\mu_\ast,b_\ast$}{u*,mu*,b*} 生成同一十次数域]\label{cor:fold-gauge-anomaly-field-equality-u-mu-b}
\leanverified{paper\_fold\_gauge\_anomaly\_field\_equality\_u\_mu\_b}
取任一满足 \(P_{10}(u_\ast)=0\) 的非平凡分歧点，并令 \(\mu_\ast\) 为相应重根，\(b_\ast\) 如 \eqref{eq:fold-gauge-anomaly-bstar-definition}。
则
\[
\boxed{\ \QQ(u_\ast)=\QQ(\mu_\ast)=\QQ(b_\ast)\ } ,
\]
且该公共数域为 \(\QQ\) 的次数 \(10\) 扩张。
\end{corollary}
\begin{proof}
由命题 \ref{prop:fold-gauge-anomaly-P10-galois-discriminant} 与命题 \ref{prop:fold-gauge-anomaly-Q10-tschirnhaus}，多项式 \(P_{10}\) 与 \(Q_{10}\) 均为次数 \(10\) 的不可约整系数多项式，故 \([\QQ(u_\ast):\QQ]=[\QQ(\mu_\ast):\QQ]=10\)。
由命题 \ref{prop:fold-gauge-anomaly-branchpoint-rational-reduction} 的有理表达 \(u_\ast=\mathcal R(\mu_\ast)\) 得 \(\QQ(u_\ast)\subseteq \QQ(\mu_\ast)\)，从而两者必相等。
另一方面由 \eqref{eq:fold-gauge-anomaly-bstar-definition} 得 \(b_\ast\in\QQ(\mu_\ast)\)；
而命题 \ref{prop:fold-gauge-anomaly-leyang-slope-square-minpoly} 给出 \(b_\ast\) 满足次数 \(10\) 的不可约整系数多项式，故 \([\QQ(b_\ast):\QQ]=10\)。
于是 \(\QQ(b_\ast)\subseteq\QQ(\mu_\ast)\) 强制 \(\QQ(b_\ast)=\QQ(\mu_\ast)\)。证毕。
\end{proof}

% --- Distillation writeback ---
\begin{proposition}[Lee--Yang 分歧定标的 Puiseux 缺口滤过与终端残数判据]\label{distill:grothendieck-cohomological_gap_filtration-folding-leyang-puiseux-terminal-residue}
沿用定理 \ref{thm:fold-gauge-anomaly-leyang-universal-kernels} 的非退化主导分歧点 \((\mu_\ast,u_\ast)\)，并取一支平方根 \(\sqrt{b_\ast}\)。令
\[
u(s):=u_\ast+s^2,
\qquad
\mu_\pm(s)=\mu_\ast\exp\bigl(\pm\sqrt{b_\ast}\,s+O(s^2)\bigr)
\]
为并合的两条 Puiseux 谱支。设一族局部审计量 \(G_m\) 在 \(s=0\) 的去心邻域内具有有限 Puiseux 缺口展开
\[
\begin{aligned}
G_m(\nu(s))=\;&a_{-1}(s^2)\frac{\mu_+(s)^m-\mu_-(s)^m}{s}
+a_0(s^2)\bigl(\mu_+(s)^m+\mu_-(s)^m\bigr)\\
&+\mu_\ast^m\sum_{\ell=1}^{L}s^\ell a_\ell(s^2)+R_m(s),
\end{aligned}
\]
其中 \(a_{-1},a_0,a_1,\ldots,a_L\) 在 \(0\) 邻域解析，且对任意紧集 \(K\subset\CC\) 有
\[
\sup_{w\in K}\left|\frac{R_m\bigl(w/(m\sqrt{b_\ast})\bigr)}{\mu_\ast^m}\right|\longrightarrow0.
\]
定义两条终端传播算子
\[
\partial_{-1}(c)(w):=2\sqrt{b_\ast}\,c\,\frac{\sinh w}{w},
\qquad
\partial_0(c)(w):=2c\cosh w .
\]
则在任意紧集 \(K\subset\CC\) 上一致成立
\[
\lim_{m\to\infty}
\frac{G_m\!\left(u_\ast+\frac{w^2}{b_\ast m^2}\right)}{m\mu_\ast^m}
=\partial_{-1}\bigl(a_{-1}(0)\bigr)(w).
\]
若进一步 \(a_{-1}(0)=0\)，则
\[
\lim_{m\to\infty}
\frac{G_m\!\left(u_\ast+\frac{w^2}{b_\ast m^2}\right)}{\mu_\ast^m}
=\partial_0\bigl(a_0(0)\bigr)(w).
\]
因此所有正 Puiseux 阶 \(s^\ell a_\ell(s^2)\) 都是 Lee--Yang 定标下的有限部伪层；它们在终端核中塌缩为零。相反，极点残数层 \(a_{-1}(0)\) 不被任何有限正阶 jet 的消失所控制，并且在 \(\mu_\ast^m\)-归一化下必须先被杀死，方可进入 \(\cosh\)-型终端层。
\end{proposition}
\begin{proof}
取 \(s=w/(m\sqrt{b_\ast})\)。由 Puiseux 展开式，紧集上一致有
\[
\frac{\mu_\pm(s)^m}{\mu_\ast^m}
=\exp\bigl(\pm w+mO(s^2)\bigr)=e^{\pm w}(1+o(1)),
\]
因为 \(m s^2=O(m^{-1})\)。于是
\[
\frac{1}{m\mu_\ast^m}\,
 a_{-1}(s^2)\frac{\mu_+(s)^m-\mu_-(s)^m}{s}
=a_{-1}(s^2)\frac{e^w-e^{-w}+o(1)}{ms}
\longrightarrow 2\sqrt{b_\ast}\,a_{-1}(0)\frac{\sinh w}{w}.
\]
同一归一化下，\(a_0\)-项为 \(O(m^{-1})\)，正阶 Puiseux 项为 \(O(m^{-1})\)，余项由假设为 \(o(m^{-1})\)，故得到第一条极限。

若 \(a_{-1}(0)=0\)，则 \(a_{-1}(s^2)=O(s^2)\)。在 \(\mu_\ast^m\)-归一化下，极点项为
\[
O(s^2)\frac{e^w-e^{-w}+o(1)}{s}=O(s)\longrightarrow0.
\]
而
\[
\frac{a_0(s^2)(\mu_+(s)^m+\mu_-(s)^m)}{\mu_\ast^m}
\longrightarrow a_0(0)(e^w+e^{-w})=2a_0(0)\cosh w.
\]
每个正阶项含因子 \(s^\ell=O(m^{-\ell})\)，故消失；余项仍由假设消失。若 \(a_{-1}(0)\ne0\)，则同一计算显示 \(\mu_\ast^m\)-归一化下极点项在一般 \(w\) 上具有 \(m\)-级增长，故不能被正阶有限 jet 抵消。命题得证。
\end{proof}

% --- End distillation writeback ---

\endinput
